\section{Spectral Model of the Zeros}
\label{sec:modele_spectral}

This section presents the analytic and operator-theoretic framework. Four
main results: closed-form expression of the amplitudes $(\kappap)$,
exact Fredholm representation of $\Rres$, canalized operator for the
intra-cascade coupling, and explicit trace formula.

\subsection{Closed-Form Expression of the Amplitudes $\kappap(s)$}
\label{sec:kappa_p}

\begin{definition}
\label{def:kappa_p}
For each prime $p$ and each $s \in \Cplx$, we set
\begin{equation}
\boxed{\;\kappap(s) \;=\; 1 \,-\, \big(1 - p^{-s}\big)\,\exp\!\big(\Ap\,p^{-s}\big)\;}
\label{eq:kappa_p}
\end{equation}
\end{definition}

\noindent By construction,
\[
1 \,-\, \kappap(s) \;=\; \big(1 - p^{-s}\big)\,\exp\!\big(\Ap\,p^{-s}\big),
\]
and the asymptotic expansion for $p \to \infty$ gives
\begin{equation}
\kappap(s) \;=\; (1 - \Ap)\,p^{-s} \,+\, O\!\big(p^{-2s}\big).
\label{eq:kappa_asymp}
\end{equation}

\paragraph{Interpretation.}
The fundamental scalar of the cascade is neither $p^{-s}$ (Euler term)
nor $\exp(\Ap p^{-s})$ (bifurcated compensation), but their gap $\kappap$.
We have the reading
\[
1 \,-\, \kappap \;=\; \text{residual survival},
\qquad
\kappap \;=\; \text{spectral amplitude of the residue}.
\]

\subsection{Exact Fredholm Representation}
\label{sec:fredholm}

Let $\DRop(s)$ be the diagonal operator on $\ell^{2}(\Primes)$ ($\Primes$
the set of primes) defined by $\DRop(s)\,e_p = \kappap(s)\,e_p$.

\begin{theoreme}[Fredholm representation of $\Rres$]
\label{thm:fredholm}
For $\Reop(s) > 1$, the operator $\DRop(s)$ is trace class, and
\begin{equation}
\boxed{\;\det\!\big(I - \DRop(s)\big)^{-1} \;=\; \Rres(s)\;.}
\label{eq:fredholm}
\end{equation}
\end{theoreme}

\begin{proof}
For $\Reop(s) > 1$, $\sum_p |\kappap(s)| < \infty$ by the asymptotic
expansion~\eqref{eq:kappa_asymp}; the diagonal operator is therefore trace
class. Direct computation gives
\begin{align*}
\det\!\big(I - \DRop(s)\big)^{-1}
&= \prod_p (1 - \kappap)^{-1}
= \prod_p (1 - p^{-s})^{-1}\,\exp(-\Ap\,p^{-s})\\
&= \zeta(s)\,\exp\!\Big(-\sum_p \Ap\,p^{-s}\Big)
= \Rres(s),
\end{align*}
using the classical Euler product for $\zeta$ and the definition of
$\Rres$ in~\eqref{eq:Rres_def}.
\end{proof}

\paragraph{Scope of the representation.}
The proof is direct algebraic: \eqref{eq:fredholm} is a
\emph{determinantal rewriting} of the definition of $\Rres$. Its value is
to provide the \emph{spectral target} that any candidate canonical PT
operator must reproduce; it is not an independent argument for the
spectral structure of $\Rres$.

\paragraph{Numerical verification.}
For $s = 1.3 + 14.134725\,i$, $P = 10^{6}$ truncated primes, the residual
error between $\Rres(s)$ computed by the Fredholm product
\eqref{eq:fredholm} and $\Rres(s)$ computed from $\zeta$ via
\texttt{mpmath} is of the order of $2.0 \cdot 10^{-12}$. At $\sigma = 1.1$
the precision drops to $\sim 1.3 \cdot 10^{-3}$, reflecting the approach
to the boundary of absolute convergence, not a defect of the
representation.

\begin{remarque}[Non-vanishing hypothesis]
\label{rem:non_annulation}
The identification $\mathrm{zeros}(\Rres) = \mathrm{zeros}(\zeta)$
used in the sequel assumes
$\zetaplus(s)\,\zetaminus(s) \neq 0$ on the zones of interest (in
particular the critical line). This hypothesis is verified
numerically on the tested zones but is not proved; it is
\textbf{conjectural} and must be treated as such
(cf.~\S\ref{sec:limites_ouvertures}).
\end{remarque}

\subsection{Canalized Operator for the Coupling $\sum_p \Ap p^{-s}$}
\label{sec:operateur_canalise}

To model the intra-cascade coupling (the sum $\sum_p \Ap p^{-s}$)
as the spectrum of a canonical operator \emph{without adjustable parameter},
we construct
\begin{equation}
\Hilb_{\mathrm{PT}}
\;=\;
L^{2}(\Sigma_{\mathrm{pers}}, du) \,\otimes\, L^{2}\!\big(\widehat{\Integers/M\Integers}\big).
\label{eq:H_PT_def}
\end{equation}
The first component is the continuous Mellin coordinate on
$\Sigma_{\mathrm{pers}}$ (persistence surface, of genus $14$ in the
PT corpus \cite{PT_GeoSpectral}); the second is discrete, on the
characters of the multiplicative group modulo $M$ (cascade by Chinese
remainder theorem, $M$ primorial product). The inter-window coupling is
constructed by regularised Bergman kernel
\begin{equation}
B_{\mathrm{BF}}(z_i, z_j)
\;=\;
\frac{\eps_i\,\eps_j}{|z_i - z_j|^{2} + \eps_i\,\eps_j},
\qquad
\eps_i \;=\; \frac{\epsstar}{\rho_i}.
\label{eq:Bergman_BF}
\end{equation}

\begin{lemme}[Geometric determination of $\epsstar$]
\label{lem:epsstar_geom}
The regularised Bergman covariance at maximal resolution
($\eps_{\mathrm{reg}} = 10^{-3}$, packets of individual prime size)
on the atomic measure $\pi_{\log}$ of the Mellin coordinate gives
\begin{equation}
\epsstar^{\mathrm{geom}} \;=\; 1.51.
\label{eq:epsstar_geom}
\end{equation}
\end{lemme}

\begin{proof}[Sketch]
The regularised Bergman covariance of the kernel~\eqref{eq:Bergman_BF} on
the atomic measure $\pi_{\log} = \sum_p \delta(u - \log p)$ is, at
resolution $\eps_{\mathrm{reg}}$, the weighted average
$\mathrm{Cov}(\eps_{\mathrm{reg}}) = \int\!\!\int B_{\mathrm{BF}}(u, u')^{2}\,
\pi_{\log}(du)\,\pi_{\log}(du')$. The numerical sweep over values of
$\eps_{\mathrm{reg}}$ by steps $10^{-3}, 10^{-2}, 10^{-1}, 1$ gives
a covariance respectively of $0.0252$, $0.0794$, $0.2408$ and
$0.680$, that is a factor $\epsstar^{\mathrm{geom}} = 1.51$ at the
maximal resolution (natural reference of the continuous model). This is
the canonical value injected into the canalized operator.
\end{proof}

\paragraph{Performance.}
With $\epsstar^{\mathrm{geom}}$ injected without adjustment, the median
error on $\sum_p \Ap\,p^{-s}$ is $0.085$ at $M = 1155$, $W = 24$ --- a
factor $\sim 85$ above the adaptive smoothed scalar baseline
($0.001$, which uses a tuned $\beta$). The contribution is not absolute
precision but the \textbf{absence of a free parameter}: \eqref{eq:epsstar_geom}
replaces an empirical tuning by a canonical geometric covariance.

\subsection{Explicit Trace Formula}
\label{sec:trace_formula}

\begin{theoreme}[Explicit trace formula]
\label{thm:trace_formula}
For all $s$ with $\Reop(s) > 1$,
\begin{equation}
\boxed{\;
-\,\dfrac{d}{ds} \log \Rres(s)
\;=\; \sum_{p} \dfrac{\kappap'(s)}{1 - \kappap(s)}
\;=\; \sum_{p} (\log p)\,\Big[\dfrac{1}{p^s - 1} \,+\, \Ap\,p^{-s}\Big]\,.\;}
\label{eq:trace_formula}
\end{equation}
\end{theoreme}

\begin{proof}
By logarithmic differentiation of
$1 - \kappap(s) = (1 - p^{-s}) \exp(\Ap p^{-s})$. Set $z = p^{-s}$,
so that $dz/ds = -\log p \cdot z$. A direct computation gives
\[
\dfrac{\kappap'(s)}{1 - \kappap(s)}
\;=\; -(\log p)\,\Big[\Ap\,z - \dfrac{z}{1 - z}\Big]
\;=\; -(\log p)\,\Big[\Ap\,p^{-s} - \dfrac{1}{p^s - 1}\Big].
\]
Summing over $p$ and changing the overall sign yields the announced
formula. A \emph{sanity check} via $-\zeta'/\zeta(s) = \sum_n
\Lambda(n)\,n^{-s} = \sum_p (\log p)/(p^s - 1)$ gives the same expression
for the first term.
\end{proof}

\paragraph{Numerical verification.}
At $s = \sigma + 14.134725\,i$ with $\sigma \in \{1.2,\,1.5,\,2,\,3\}$ and
$2\,000$ truncated primes, the discrepancy between the two forms
of~\eqref{eq:trace_formula} is of the order of $10^{-26}$ (precision
\texttt{mpmath} 25 digits). The identity is therefore exact in the
algebraic sense.

\paragraph{Berry--Keating reading.}
The dilation operator
$\HPTBK = (u\,p_u + p_u\,u)/2 = -i(u\partial_u + 1/2)$ generates the flow
$e^{-it\HPTBK} f(u) = e^{-t/2}\,f(e^{-t}u)$. On the atomic measure
$\nu = \sum_p \delta(u - \log p)$, the periodic orbits of the flow
starting from $\log p_0$ have lengths $\ell_p = k \log p_0$, $k \ge 1$,
corresponding to the powers $p_0^{k}$. The term
$(\log p)/(p^s - 1) = \sum_{k \ge 1} (\log p)\,p^{-ks}$ identifies then
with the sum over these orbits, and $\Ap\,p^{-s}$ represents the
$\qplus/\qminus$ holonomic coupling contribution per prime.

\subsection{Regularised Operator Trace}
\label{sec:trace_reg}

The compact writing of the trace formula~\eqref{eq:trace_formula} in
terms of the resolvent of $\HPTBK$ and of the operator $\DRop$ uses an
atomic Mellin projection
\begin{equation}
\Pi : L^{2}([\umin, \umax], du) \,\to\, \ell^{2}(\Primes_\gamma),
\qquad \Pi(f)_p \;=\; f(\log p)\,\sqrt{\log p},
\label{eq:Pi_def}
\end{equation}
where $\Primes_\gamma = \{p \text{ prime} : \umin \le \log p \le
\umax(\gamma)\}$. The factor $\sqrt{\log p}$ is the square root of the
atomic multiplicative Haar measure. The regularised trace then writes
\begin{equation}
\boxed{\;
\Treg(s) \;:=\; \Tr_{\mathrm{reg}}\!\Big[(s - i\HPTBK)^{-1}\,\DRop(s)\Big]
\;=\; \sum_{p \in \Primes_\gamma} \kappap(s)\,(\log p)\,\Big[\dfrac{1}{p^s - 1} + \Ap\,p^{-s}\Big]\,.\;}
\label{eq:Treg}
\end{equation}
In $\Reop(s) > 1$, $\Treg(s) \equiv -d \log \Rres/ds$. On the critical
line $\Reop(s) = 1/2$, the prime sum is not absolutely convergent; it is
however distributionally convergent as an element of
$\mathcal{S}'(\Reals_t)$, and the Hadamard regularisation gives a
functional sense to $\Treg(1/2 + it)$.

\paragraph{Three equivalent regularisations.}
We will use three regularisations interchangeable in the distributional
sense.
\begin{description}[nosep,leftmargin=*]
\item[($R_1$) Smearing by test function.] For $\phi_\eps$ a compactly
supported approximation of the identity with width $\eps$, we set
$I_\eps(\gamma_0) := \int \phi_\eps(t - \gamma_0)\,F(1/2 + it)\,dt$.
The peaks $I_\eps(\gamma_n)$ behave as $c/\eps$ when $\eps \to 0^{+}$
(smearing regularisation).
\item[($R_2$) Abel-Hadamard summation.] With weight $p^{-\eps}$, $\eps > 0$,
$F_\eps(t) := \sum_p (\log p)\,p^{-\eps}\,[1/(p^{1/2+it} - 1) + \Ap\,p^{-1/2-it}]$.
By Fourier transform, $F_\eps$ is the convolution of $F$ by the
Gaussian $\phi_\eps$: $R_2$ is equivalent to $R_1$ modulo the choice
of regularising kernel.
\item[($R_3$) Subtraction of the principal pole.] $\widetilde F(s) := F(s) - 1/(s - 1)$
isolates the simple pole at $s = 1$ and applies the Hadamard extraction
to the remaining part.
\end{description}

The three regularisations coincide in the distributional sense at the
poles of $\Rres$. The strict functional proof of the analytic
continuation of $\Treg$ to the critical line --- beyond distributional
convergence --- remains open (\S\ref{sec:limites_ouvertures}).

\paragraph{Self-correction: G\_HP4.d falsified.}
A first test posed the naive identity ``direct spectrum of $\HPTBK$
bounded by double barrier $= \{\gamma_n\}$''. This test was
falsified: the direct spectrum is regular (Bohr--Sommerfeld
asymptotically linear), while the $\gamma_n$ are irregular, and
the ratio between the two spectra diverges. This falsification motivated
the move to the distributional regularised trace: the zeros of
$\zeta$ are \emph{distributional poles} of $\Treg(s)$, not direct
eigenvalues of $\HPTBK$. The falsifiability discipline of the
programme includes this correction on the same footing as the
Dirichlet falsification of \S\ref{sec:auto_correction_dirichlet}.

